\chapter{Conclusion}
NimbusVault implements a practical connection between secure file management and evidence-grounded document intelligence. PostgreSQL and Express control identity, organisation and sharing; private S3 retains original bytes; existing event-driven AWS services record bounded processing stages. The private MiniLM/SQLite engine provides dense, lexical and hybrid discovery, while Qwen performs genuine local generation with source-attributed answers and visible failure modes.

The strongest result is not an isolated language-model response but an integrated boundary: readable documents determine retrieval scope, citations are mapped to currently authorised files, and conversations retain ownership checks. Evaluation provides concrete evidence for this integration. Separately executed suites recorded 460 passes, browser verification covered 54 route/viewport combinations, hybrid retrieval accepted all 69 relevant top-1 cases and rejected fifteen negatives, and the final forty-case CPU AI run met its strict generated/abstention criteria. These are bounded regression results, not universal correctness claims.

Standalone repair and native Linux preparation also improved reproducibility. Copied assets passed integrity checks, recovery data remained intact, and active-generation shutdown and restart were demonstrated locally. Historical AWS rehearsal evidence remains distinct from read-only configuration observations. Full EC2 hosting is still awaiting explicit authorisation and acceptance; this report does not convert the proposed host into an accomplished deployment.

The project demonstrates applied cloud computing through its separation of storage, processing, identity and inference, and applied AI through measured retrieval and grounded response handling. Its next steps follow directly from the evidence: finish the approved cloud acceptance, expand independent evaluation and address residual security and capacity limits without abandoning preservation or financial constraints.
\FloatBarrier
